\SourceLeaf{244}%
perficie data, cuius indoles exprimitur per aequationem \(P\mathrm{d}x+Q\mathrm{d}y+R\mathrm{d}z=0\), etiam variationes \(\delta x,\delta y,\delta z\) satisfacere debent aequationi \(P\delta x+Q\delta y+R\delta z=0\), unde per principia nota facile colligitur, differentialia
\[
\mathrm{d}\frac{\mathrm{d}x}{\BareRoot{(\mathrm{d}x^2+\mathrm{d}y^2+\mathrm{d}z^2)}},\quad
\mathrm{d}\frac{\mathrm{d}y}{\BareRoot{(\mathrm{d}x^2+\mathrm{d}y^2+\mathrm{d}z^2)}},\quad
\mathrm{d}\frac{\mathrm{d}z}{\BareRoot{(\mathrm{d}x^2+\mathrm{d}y^2+\mathrm{d}z^2)}}
\]
resp. quantitatibus \(P,Q,R\) proportionalia esse debere. Iam sit \(\mathrm{d}r\) elementum lineae curvae, \(\lambda\) punctum in superficie sphaerica repraesentans directionem huius elementi, \(L\) punctum in superficie sphaerica repraesentans directionem normalis in superficiem curvam; denique sint \(\xi,\eta,\zeta\) coordinatae puncti \(\lambda\), atque \(X,Y,Z\) coordinatae puncti \(L\) respectu centri sphaerae. Ita erit
\[
\mathrm{d}x=\xi\mathrm{d}r,\quad\mathrm{d}y=\eta\mathrm{d}r,\quad\mathrm{d}z=\zeta\mathrm{d}r
\]
unde colligimus, differentialia illa fieri \(\mathrm{d}\xi,\mathrm{d}\eta,\mathrm{d}\zeta\). Et quum quantitates \(P\mathord{.}\ Q\mathord{.}\ R\) proportionales sint ipsis \(X,Y,Z\), character lineae brevissimae consistit in aequationibus
\[
\frac{\mathrm{d}\xi}{X}=\frac{\mathrm{d}\eta}{Y}=\frac{\mathrm{d}\zeta}{Z}
\]
Ceterum facile perspicitur, \(\BareRoot{(\mathrm{d}\xi^2+\mathrm{d}\eta^2+\mathrm{d}\zeta^2)}\) aequari arculo in superficie sphaerica, qui mensurat angulum inter directiones tangentium in initio et fine elementi \(\mathrm{d}r\), adeoque esse \(=\frac{\mathrm{d}r}{\rho}\), si \(\rho\) denotet radium curvaturae in hoc loco curvae brevissimae; ita fiet
\[
\rho\mathrm{d}\xi=X\mathrm{d}r,\quad\rho\mathrm{d}\eta=Y\mathrm{d}r,\quad\rho\mathrm{d}\zeta=Z\mathrm{d}r
\]
\Article{15}
Supponamus, in superficie curva a puncto dato \(A\) proficisci innumeras curvas brevissimas, quas inter se distinguemus per angulum, quem constituit singularum elementum primum cum elemento primo unius ex his lineis pro prima assumtae: sit \(\varphi\) ille angulus, vel generalius functio illius anguli, nec non \(r\) longitudo talis lineae brevissimae a puncto \(A\) usque ad punctum, cuius coordinatae sunt \(x,y,z\). Quum itaque valoribus determinatis variabilium \(r,\varphi\) respondeant puncta determinata superficiei, coordinatae \(x,y,z\) considerari possunt tamquam functiones ipsarum \(r,\varphi\). Notationes \(\lambda,L,\xi,\eta,\zeta,X,Y,Z\) in eadem significa%
